% SPDX-License-Identifier: CC-BY-4.0
% Copyright (c) 2026 Martin Schröder <info@swedishembedded.com>

\section{Training}\label{sec:training}

Splinter chooses the objective and the data; Brain runs the optimisation. LoRA
supervised fine-tuning is Brain's \code{ChatFineTune} and DPO is its
\code{PreferenceFineTune}; both are called from the model adapter crate.\cl{T01} The
defaults are a learning rate of $1\times10^{-4}$, rank 8, alpha twice the rank and no
weight decay;\cl{T02} the step budget is three passes over the records clamped to between
40 and 2000 steps.\cl{T03} A candidate \emph{continues the champion}: its training starts
from the parent release's adapter,\cl{T06} mixes in a replay of each earlier release's
records (a quarter of each earlier release's records, entering a quarter of training draws),\cl{T04} and mixes in a rehearsal of the base's own
answers (a quarter of draws, half of it code-graded form tasks).\cl{T05} A monitored run
keeps the best evaluation's adapter.\cl{T07}

\paragraph{Consolidation.} A policy that keeps learning would grow its adapter, or must
retrain from scratch on every claim. Splinter therefore consolidates a generation into
the weights in one of two modes: a full fine-tune of every weight of the generation's
base, or merging the adapter into the base.\cl{T08} The full mode defaults to a learning
rate of $1\times10^{-5}$ for six epochs with a retention replay share of a quarter;\cl{T09} a
generation is consolidated automatically after more than 2000 trained records.\cl{T10}
Both are carried out by Brain (\code{ChatFullFineTune}, \code{merge\_adapter}).\cl{T11}

\section{Release gate, releases, rollback and lineage}\label{sec:gate}

\paragraph{The gate.} A candidate is released only if four checks all pass, each over
material it was not trained on.\cl{G01} A check that was not measured fails with the
reason \emph{not measured}.\cl{G02}
\emph{Improvement}: a one-sided exact paired sign test at $\alpha=0.05$ over discordant
clusters, a family of items counting once; the candidate must also have solved at least as
many of the suite's assigned tasks as its champion, an attempt withheld as exposed or as a
harness failure counting as unsolved for the arm that produced it.\cl{G03}
\emph{Retention}: accuracy on every earlier release's held-out suite may not drop by more
than 0.05.\cl{G04}
\emph{Anchor}: accuracy on the frozen suite of general tasks may not drop by more than
0.02; the report carries a 95\% bootstrap interval over 4000 resamples.\cl{G05}
\emph{Serve}: plain \code{brain serve} must report the candidate's adapter digest and
agree with the in-process answers on at least 75\% of up to eight re-answered held-out
tasks.\cl{G06} Repository exams use the same four checks, with the exam families as the
improvement suite and agreement judged on the graded verdicts.\cl{G07} The bootstrap seed
is drawn when not given and recorded;\cl{G08} a failed retention check queues the
forgotten concepts;\cl{G09} and only a candidate trained from the current champion can
replace it.\cl{G10}

\paragraph{Releases.} A release manifest (format \code{splinter-release-v4}) records the
artifact, parent, candidate, datasets, replay, rehearsal, training settings, the whole
gate report, metrics, terms, provenance and generation.\cl{X01} A release is written
once and an existing one is refused;\cl{X02} unrestricted distribution is refused unless
the terms of everything it was made from permit it.\cl{X03} The default policy is the alias
\code{policy:default}; moving an alias is a compare-and-set from the release the gate
measured against.\cl{X04} After every move the serving manifest
\code{<state>/serving/<alias>.json} is rewritten, so \code{brain serve --adapter-manifest}
follows the release without a restart.\cl{X05} \code{splinter rollback} moves the alias to
the release the current one was trained from and refuses when there is none.\cl{X06}
\code{splinter lineage} walks any artifact up to its source bytes and down to everything
derived from it, computed from the stores on every call.\cl{X07}
